% !TEX TS-program = xelatex
% !BIB TS-program = bibtex
% Q1 checkpoint, DLV-05 (W1). Source of truth: cc-research docs/paper/q1/. Numbers only through \nm{key} (numbers.json);
% lint_report.py rejects any other numeral. Prior-work sentences carry \cite or end with \ours (our own statement).
\documentclass[12pt,letterpaper]{article}
\usepackage{style/dsc180reportstyle} % import dsc180reportstyle.sty

% Generated by build_numbers.py from numbers.json. Do not edit.
% cc-research HEAD aef649241abed84577af7443c8481aab5420b9fd
\expandafter\def\csname nm@cc.arms\endcsname{30}
\expandafter\def\csname nm@cc.completed\endcsname{22}
\expandafter\def\csname nm@cc.s0-adj\endcsname{0.584}
\expandafter\def\csname nm@cc.final-adj\endcsname{0.321}
\expandafter\def\csname nm@cc.usd\endcsname{\$12.57}
\expandafter\def\csname nm@cc.cyc.entered\endcsname{8}
\expandafter\def\csname nm@cc.cyc.keyerr\endcsname{8}
\expandafter\def\csname nm@bnd08.runs\endcsname{30}
\expandafter\def\csname nm@bnd08.m3.labels\endcsname{81}
\expandafter\def\csname nm@bnd08.m3.data\endcsname{20}
\expandafter\def\csname nm@bnd08.m3.n\endcsname{101}
\expandafter\def\csname nm@bnd08.model\endcsname{gpt-4o-2024-08-06}
\expandafter\def\csname nm@ab06.agents\endcsname{4}
\expandafter\def\csname nm@ab06.cc-entered\endcsname{1}
\expandafter\def\csname nm@ab06.systems\endcsname{5}
\expandafter\def\csname nm@ab06.p9.entered\endcsname{4}
\expandafter\def\csname nm@ab06.p9.losing\endcsname{3}
\expandafter\def\csname nm@ab06.p9.min-scen\endcsname{5}
\expandafter\def\csname nm@ab06.p9.threshold\endcsname{0.10}
\expandafter\def\csname nm@ab06.causaliq.drop\endcsname{0.251}
\expandafter\def\csname nm@ab06.cc.drop\endcsname{0.167}
\expandafter\def\csname nm@ab06.matmcd.drop\endcsname{0.122}
\expandafter\def\csname nm@ab06.vashishtha.drop\endcsname{0.020}
\expandafter\def\csname nm@ab06.matmcd.dv1\endcsname{0.6855}
\expandafter\def\csname nm@ab06.matmcd.dv1-data\endcsname{0.7199}
\expandafter\def\csname nm@ab01.pc.published\endcsname{0.333}
\expandafter\def\csname nm@ab01.pcmaxp.controlled\endcsname{0.463}
\expandafter\def\csname nm@ab01.fci.authors\endcsname{0.578}
\expandafter\def\csname nm@ab01.fci.expect\endcsname{0.357}
\expandafter\def\csname nm@ab01.fci.gain\endcsname{0.221}
\expandafter\def\csname nm@ab01.files.semantic\endcsname{250}
\expandafter\def\csname nm@ab01.files.formula\endcsname{250}
\expandafter\def\csname nm@ab01.files.synthetic\endcsname{100}
\expandafter\def\csname nm@ben13.reproduced\endcsname{4}
\expandafter\def\csname nm@ben13.cells\endcsname{4}
\expandafter\def\csname nm@ben13.for.random-level\endcsname{3}
\expandafter\def\csname nm@ben13.echo.min\endcsname{0.90}
\expandafter\def\csname nm@ben13.echo.max\endcsname{0.95}
\expandafter\def\csname nm@s2.families\endcsname{10}
\expandafter\def\csname nm@s2.tables\endcsname{30}
\expandafter\def\csname nm@s2.arms-per\endcsname{3}
\expandafter\def\csname nm@s2.rows\endcsname{1,000}
\expandafter\def\csname nm@s2.vars.min\endcsname{20}
\expandafter\def\csname nm@s2.vars.max\endcsname{25}
\expandafter\def\csname nm@s2.edges.min\endcsname{35}
\expandafter\def\csname nm@s2.edges.max\endcsname{58}
\expandafter\def\csname nm@s1.sachs.rows\endcsname{853}
\expandafter\def\csname nm@s1.sachs.vars\endcsname{11}
\expandafter\def\csname nm@s1.sachs.edges\endcsname{19}
\expandafter\def\csname nm@s1.sachs.arms\endcsname{4}
\expandafter\def\csname nm@s1.sachs.bnlearn\endcsname{17}
\expandafter\def\csname nm@s1.sachs.shared\endcsname{16}
\expandafter\def\csname nm@s1.sachs.reversed\endcsname{1}
\expandafter\def\csname nm@s1.sachs.arena\endcsname{20}
\expandafter\def\csname nm@s1.abalone.rows\endcsname{4,177}
\expandafter\def\csname nm@s1.abalone.vars\endcsname{8}
\expandafter\def\csname nm@s1.abalone.edges\endcsname{7}
\expandafter\def\csname nm@s1.autompg.rows\endcsname{392}
\expandafter\def\csname nm@s1.autompg.vars\endcsname{5}
\expandafter\def\csname nm@s1.autompg.edges\endcsname{5}
\expandafter\def\csname nm@s1.concrete.rows\endcsname{1,030}
\expandafter\def\csname nm@s1.concrete.vars\endcsname{9}
\expandafter\def\csname nm@s1.concrete.edges\endcsname{8}
\expandafter\def\csname nm@s1.dwd.rows\endcsname{349}
\expandafter\def\csname nm@s1.dwd.vars\endcsname{5}
\expandafter\def\csname nm@s1.dwd.edges\endcsname{4}
\expandafter\def\csname nm@s1.autompg.arena.vars\endcsname{4}
\expandafter\def\csname nm@s1.autompg.arena.edges\endcsname{3}
\expandafter\def\csname nm@s1.autompg.arena.rows\endcsname{398}
\expandafter\def\csname nm@s1.named.tables\endcsname{12}
\expandafter\def\csname nm@s1.named\endcsname{4}
\expandafter\def\csname nm@s1.named.arms-per\endcsname{3}
\expandafter\def\csname nm@cb.k562.cells\endcsname{162,751}
\expandafter\def\csname nm@cb.k562.genes\endcsname{622}
\expandafter\def\csname nm@cb.k562.controls\endcsname{10,691}
\expandafter\def\csname nm@cb.rpe1.cells\endcsname{162,733}
\expandafter\def\csname nm@cb.rpe1.genes\endcsname{383}
\expandafter\def\csname nm@cb.rpe1.controls\endcsname{11,485}
\expandafter\def\csname nm@cc.head\endcsname{aef649241abe}

\makeatletter
\newcommand{\nm}[1]{\@ifundefined{nm@#1}{\PackageError{numbers}{undefined number key #1}{}}{\csname nm@#1\endcsname}}
\makeatother
\newcommand{\ours}{}
\setlength{\emergencystretch}{1em}
\makeatletter
\widowpenalty=\@M\relax
\clubpenalty=\@M\relax
\makeatother

%%%%%%%%%%%%%%%%%%%%%%%%%%%%%%%%%%%%%%%%%%%%%%%%%%%%%%%%
%%%% Title and Authors
%%%%%%%%%%%%%%%%%%%%%%%%%%%%%%%%%%%%%%%%%%%%%%%%%%%%%%%%

\title{Crucis: An Independent Framework\\ for Auditing Causal-Discovery Systems}

\author{Uzair Gheewala \\
  {\tt \studentemail} \\}
\newcommand{\studentemail}{ugheewala@ucsd.edu}

\newcommand{\checkpointbibstyle}{\small\linespread{0.96}\selectfont}
\definecolor{StageHighlight}{gray}{.9}
\begin{document}
\maketitle

%%%%%%%%%%%%%%%%%%%%%%%%%%%%%%%%%%%%%%%%%%%%%%%%%%%%%%%%
%%%% Abstract and Links
%%%%%%%%%%%%%%%%%%%%%%%%%%%%%%%%%%%%%%%%%%%%%%%%%%%%%%%%

\begin{abstract}
A causal discovery system may use numerical patterns, variable names, or both to decide
which arrows to draw. A final benchmark score tells us little about how it reached those
decisions. We introduce Crucis, an independent framework for auditing externally developed
systems by changing one source of information at a time where the comparison permits,
and following changes through recorded stages. The contracts for connecting systems and
defining experiments are implemented and tested. Recorded audits illustrate what the
approach can reveal. On the CausalArena semantic cases we examined, using column order
to settle uncertain arrow directions raised the Fast Causal Inference algorithm's
score on identical graphs. Changing
variable names while keeping the data fixed lowered direction scores in several recorded
comparisons. These used different measurement setups, and the Causal-Copilot result was
exploratory. In a separate analysis of completed Causal-Copilot runs, the largest paired
drop in the score for which variables were connected occurred during bootstrap refinement,
before language-model review.
\end{abstract}

\begin{center}
Framework: \url{https://github.com/crucis-lab/crucis}\\
Research evidence: \url{https://github.com/uzairgheewala/cc-research}
\end{center}

\maketoc
\clearpage

%%%%%%%%%%%%%%%%%%%%%%%%%%%%%%%%%%%%%%%%%%%%%%%%%%%%%%%%
%%%% Main Contents
%%%%%%%%%%%%%%%%%%%%%%%%%%%%%%%%%%%%%%%%%%%%%%%%%%%%%%%%

\section{Introduction}
\label{sec:intro}

If we change one variable, what else moves? When an experiment is impractical,
causal discovery can propose a graph from the available data. Language models can
now participate in building that graph. Causal-Copilot uses one to inspect data,
help select and configure an algorithm, review the graph, and report the result
\citep{wang2025causalcopilot}. Other systems use a model to propose constraints,
a causal order, or edits to a data-driven graph
\citep{khatibi2024alcm,shen2024matmcd,vashishtha2025causal}.

\subsection{What should an audit tell us?}
The same score can reflect patterns in the data or reliance on clues in the input. Variable scales can reveal causal order
in simulated benchmarks \citep{reisach2021beware}. Column order can affect a search
algorithm's output \citep{kitson2024impact}, and language models can reconstruct
familiar graphs from names alone \citep{srivastava2025llm}. An independent audit
should make these sources distinguishable. Crucis is designed to do that across
externally developed systems. Causal-Copilot is the worked example for tracing
changes inside one system.

\subsection{How does Crucis examine a result?}
A connection records what a system receives and returns. An experiment then states
what changes, what stays fixed, and how outputs will be compared. If intermediate
records are available, the audit can locate where a change appeared. Each level
answers a different question.

\begin{description}
  \item[L0 \enspace Score the graph.] Compare the output with the stated reference.
  \item[L1 \enspace Change the input clues.] Keep the numbers fixed while changing
    names, column order, or file paths under the specified controls.
  \item[L2 \enspace Trace the changes.] Use available stage records to locate added,
    removed, or reversed links. Link them to logged operations where the record permits.
  \item[L3 \enspace Trace claims.] The design calls for tracing individual claims
    back to their sources.
  \item[L4 \enspace Search for failures.] The design calls for adversarial tests of
    where a system's claims break down.
\end{description}

The system connection and experiment definition contracts are implemented and
tested. Execution and collection remain in design, as do the L3 and L4 capabilities
described here. The recorded audits exercise parts of the design on retained
evidence. Native execution across systems, validation of the auditor, and the
broader multi-system study remain ahead.

An audit reports what information the system had, such as data alone or data
with interventions. It also reports what the comparison can establish. A graph
used to simulate data can support an accuracy comparison. When published
reference graphs disagree, the score measures agreement with the specified
version. Section~\ref{sec:data} explains which reference each audit uses.

\subsection{What have the recorded audits shown?}
The audits cover \nm{ab06.systems} language-model systems and the CausalArena
and CausalBench benchmarks. Each example below answers a different question.

\begin{description}
  \item[Can the scoring rule change the result?] On the evaluated CausalArena
    semantic subset, using column order to settle uncertain directions raised the directed F1 score of Fast
    Causal Inference (FCI) by \nm{ab01.fci.gain} relative to averaging over those directions
    on the same graphs. See Appendix~\ref{app:scoring} for the scoring comparison.
  \item[Can the names change the result?] With numerical data fixed, swapping names
    lowered directed F1 by at least \nm{ab06.p9.threshold} in \nm{ab06.p9.losing} of
    \nm{ab06.p9.entered} eligible entries. The entries use different cohorts,
    scoring rules, and output stages, including an earlier exploratory
    Causal-Copilot comparison. Appendix~\ref{app:names} explains how to read the tally.
  \item[Where did the graph change?] In the completed Causal-Copilot main-control
    runs, the largest paired loss in adjacency F1 occurred during bootstrap
    refinement, before language-model review. Appendix~\ref{app:stages} follows
    the relevant stages and accounts for crashes.
\end{description}

These results show why the pieces of an audit belong together. The scoring example
depends on holding the graphs fixed. The name comparison depends on holding the data
fixed and knowing which output was measured. The bootstrap result depends on
following the same completed runs through their intermediate graphs. A final score
alone cannot make those distinctions.

\section{Prior work}
\label{sec:prior}

\subsection{How causal discovery is evaluated, and what benchmarks reward}

Greedy Equivalence Search compares classes of graphs using a score and can identify
the generating structure under the conditions of its large-sample result
\citep{chickering2002optimal}.
Benchmarks with known graphs provide a way to check recovery, but the inputs
can reveal features of the reference graph \citep{reisach2021beware,kitson2024impact}. In common simulated settings the variance of the variables grows along the causal order, so sorting by variance matches
continuous-optimisation methods \citep{reisach2021beware}, and a scale-invariant version of the same artefact is also high in common settings
and can be exploited \citep{reisach2023scale}. The order of the columns in the input data alone changes the graph that hill-climbing search
learns, typically more than a change in sample size does \citep{kitson2024impact}. Real-data evaluation leans on very few datasets, above all the Sachs
protein-signalling data, whose reference graph is not definitive \citep{brouillard2025landscape,sachs2005causal}.
CausalArena collects synthetic, semantic and formula-grounded models with a public leaderboard
\citep{li2026causalarena}, and CausalBench scores gene-network inference on held-out perturbation data
\citep{chevalley2025causalbench,replogle2022mapping}.

\subsection{Language models as a source of causal knowledge}

Language models answer pairwise causal questions about named variables well, and part of that skill is memorised
\citep{kiciman2023causal}; they may repeat causal statements without performing causal inference
\citep{zecevic2023causal}. When a model is treated as an imperfect expert, the guarantee that the true graph is kept holds
only under an assumed error model, and in the authors' own case study the model was poorly calibrated and the true
graph was excluded on one benchmark \citep{long2023causal}. Asking for a causal order instead of individual edges avoids
confusing direct with indirect effects \citep{vashishtha2025causal}. Controls for memorisation now exist: neutral aliases in
the model's queries \citep{kitson2026llm}, graphs published after a model's training cut-off \citep{srivastava2025llm},
random graphs carrying real concept names \citep{li2026llm}, and names moved onto the wrong columns with the data
unchanged, under which models' accuracy collapses \citep{dong2025care}. Our audit battery uses wrong-name reassignment alongside separate layout controls where the recorded evidence
supports those comparisons\ours.

\subsection{Causal agents and how they are evaluated}

Causal-Copilot automates the whole analysis and states that the model's recommendations are soft suggestions that may not
override the high-confidence decisions of its bootstrap procedure \citep{wang2025causalcopilot}. ALCM, MATMCD and Causal Modelling Agents let a model refine or
constrain a data-driven graph \citep{khatibi2024alcm,shen2024matmcd,abdulaal2024causal}, and open libraries expose
model-based causal suggestions \citep{causaliq_knowledge,pywhy_llm}. CausalSteward evaluates Causal-Copilot as a baseline \citep{tagliapietra2026causalsteward}. Repeated runs of
the same agent can disagree, so single-run comparisons can misrank systems \citep{mehta2026agents}. A position paper
argues that agents should not supply edges or causal conclusions at all and asks for audit trails
\citep{zheng2026causal}.

\paragraph{Related benchmarks and auditing frameworks.}
Heuresis's PC-centred trace decomposition offers a precedent for L2, alongside
instrument checks and a name-only leakage probe in its interactive evaluation
on anonymous, randomly ordered simulations \citep{ma2026heuresis}.
ACDB separates observational recovery, DAG recovery and intervention efficiency in a benchmark
with budgeted interventions, and tests how graph density affects a structure-blind
baseline \citep{deb2026acdb}.
TCD-Arena provides a modular test kit for robustness to controlled time-series assumption violations
\citep{stein2026tcdarena}. Causal-Audit uses time-series assumption diagnostics and calibrated risk to recommend a
method or abstain before discovery within its evaluated decision pipeline \citep{ruiz2026causalaudit}.
The design question for Crucis is how to use the records an external system can
supply to make a stage comparison, and how to report where those records leave
the comparison unresolved\ours.

\subsection{How much an edge can be trusted}

Bootstrap resampling gives confidence measures for features of a learned network \citep{friedman1999data}, but the
resulting edge frequencies can be miscalibrated \citep{kummerfeld2019simulations,banerjee2025calibrated}. Language
models asked about direct edges give high stated confidence to most of their false-positive edges; agreement across
prompts is better calibrated, though not significantly so \citep{kumar2026plausibility}. An audit therefore has to test whether a system's confidence means anything, not assume
it\ours.

\subsection{Auditing practice the protocol borrows}

Metamorphic testing checks a program that has no oracle for each output by checking relations between the outputs of
related inputs \citep{chen1998metamorphic,segura2016survey}; renaming or reordering columns while holding the numbers
fixed is such a relation\ours. Behavioural testing of language-processing models crosses capabilities with test types that
include invariance tests, and finds failures that held-out accuracy hides \citep{ribeiro2020beyond}. Model cards and datasheets document a model's
evaluation and a dataset's origin \citep{mitchell2019model,gebru2021datasheets}, internal algorithmic audits define
staged audit artefacts \citep{raji2020closing}, and PROV-DM records which agent and activity produced a piece of data
\citep{moreau2013prov}. Our certificate combines these ideas for causal claims: system identity, battery hash, every
number with its source, and the command that reruns it\ours.

\section{Data}
\label{sec:data}

The audit draws on a suite built from CausalArena's semantic scenarios, Sachs and
\nm{s1.named} other named datasets that ship with Causal-Copilot, and CausalBench's perturbation data. Table~\ref{tab:data}
shows what each audit uses as a reference, from a known simulated graph to held-out interventions. Shapes come from suite manifests and
retained CausalBench filtering metadata (Appendix~\ref{app:numbers}). The checkpoint binds these sources by hash.

\begin{table}[htbp]
\caption{Selected inputs by audit. The semantic rows describe the Causal-Copilot name battery. For CausalBench, cells
and perturbed genes are counted after filtering; held-out interventions provide an evaluation regime without a
reference graph. Other tiers describe graph reliability (Section~\ref{sec:intro}).}
\label{tab:data}
\centering
{\footnotesize\setlength{\tabcolsep}{3pt}
\begin{tabular}{llrrrll}
\toprule
Group & Dataset & Rows & Variables & Ref.\ edges & Tier & Licence \\
\midrule
CausalArena & \nm{s2.families} semantic scenarios & \nm{s2.rows} & \nm{s2.vars.min}--\nm{s2.vars.max} & \nm{s2.edges.min}--\nm{s2.edges.max} & exact & CC BY 4.0 \\
\midrule
Named & Sachs & \nm{s1.sachs.rows} & \nm{s1.sachs.vars} & \nm{s1.sachs.edges} & contested & none stated \\
 & Abalone & \nm{s1.abalone.rows} & \nm{s1.abalone.vars} & \nm{s1.abalone.edges} & provenance & CC BY 4.0 \\
 & Auto MPG & \nm{s1.autompg.rows} & \nm{s1.autompg.vars} & \nm{s1.autompg.edges} & mismatched & CC BY 4.0 \\
 & Concrete strength & \nm{s1.concrete.rows} & \nm{s1.concrete.vars} & \nm{s1.concrete.edges} & provenance & CC BY 4.0 \\
 & DWD climate & \nm{s1.dwd.rows} & \nm{s1.dwd.vars} & \nm{s1.dwd.edges} & provenance & see text \\
\midrule
CausalBench & K562 & \nm{cb.k562.cells} & \nm{cb.k562.genes} & -- & held-out interventions & CC BY 4.0 \\
 & RPE1 & \nm{cb.rpe1.cells} & \nm{cb.rpe1.genes} & -- & held-out interventions & CC BY 4.0 \\
\bottomrule
\end{tabular}}
\end{table}

\paragraph{CausalArena semantic suite.} CausalArena releases structural causal models whose variables carry names from
a written scenario, with observational exports \citep{li2026causalarena}. We take the \nm{s2.families} semantic
scenarios of our test split, the first replicate of the observational-only release, \nm{s2.rows} rows each. Each scenario is
written \nm{s2.arms-per} times with byte-identical numbers: under its real names, under neutral names, and under its
own names moved onto the wrong columns by a derangement chosen to reduce pairs adjacent under both readings
(\nm{s2.tables} tables). Because the released column order is a causal order, the columns are shuffled with a fixed
seed, and the file paths are opaque codes. The reference is the graph the data were simulated from, so scores are
accuracy. The layout audit of the public leaderboard uses CausalArena's public files instead (\nm{ab01.files.semantic}
semantic, \nm{ab01.files.formula} formula and \nm{ab01.files.synthetic} synthetic). The dataset is released under
CC BY 4.0 on Hugging Face (\texttt{LAMDA-Tabular/CausalArena}).

\paragraph{Named datasets.} Causal-Copilot ships Sachs and \nm{s1.named} other tables with reference graphs, and we use
its copies \citep{wang2025causalcopilot}. Sachs \citep{sachs2005causal} has \nm{s1.sachs.arms} arms (the original, a copy under
another path, neutral names, and moved names). Its reference is contested: the repository's graph has
\nm{s1.sachs.edges} edges, the bnlearn version \nm{s1.sachs.bnlearn}, sharing \nm{s1.sachs.shared} adjacencies with
\nm{s1.sachs.reversed} reversed, and CausalArena uses a version with \nm{s1.sachs.arena} edges, so a score on Sachs is
agreement with one version, not accuracy. Abalone \citep{nash1994abalone}, Auto MPG \citep{quinlan1993autompg} and
Concrete compressive strength \citep{yeh1998concrete} come from the UCI repository under CC BY 4.0, and the climate
table matches the Deutscher Wetterdienst pairs of the T\"ubingen cause-effect benchmark \citep{mooij2016distinguishing}.
The DWD open-data terms are CC BY 4.0; the copy's lineage is inferred from its match to the benchmark pairs. Each has \nm{s1.named.arms-per} arms (real, neutral, swapped). Their
reference graphs are expert graphs with nothing to compare against (provenance tier), except Auto MPG, whose reference
differs from CausalArena's version (\nm{s1.autompg.arena.vars} variables, \nm{s1.autompg.arena.edges} edges,
\nm{s1.autompg.arena.rows} rows), so its scores must not be set beside CausalArena's row. The checked sources did not identify a licence for the Sachs data.

\paragraph{CausalBench.} CausalBench uses the CRISPRi Perturb-seq screens of \citet{replogle2022mapping} in the
K562 and RPE1 cell lines and scores a predicted network on held-out interventions by a mean Wasserstein distance
and a false omission rate \citep{chevalley2025causalbench}. In the retained metadata for CausalBench's filtering, the
K562 data have \nm{cb.k562.cells} cells (\nm{cb.k562.controls} unperturbed controls) and \nm{cb.k562.genes} perturbed
genes, and RPE1 \nm{cb.rpe1.cells} cells (\nm{cb.rpe1.controls} controls) and \nm{cb.rpe1.genes} genes. The screens are
released under CC BY 4.0 and the benchmark code under Apache 2.0. Our audit of this benchmark is cited from its result
file. All \nm{ben13.reproduced} published MeanDifference rows match within the recorded tolerances.
At the smaller evaluated edge budget, MeanDifference reaches \nm{ben13.echo.min}--\nm{ben13.echo.max} of a
test-informed reference that ranks pairs by knockdown effects measured on the test split. Its false omission
rate stays within the audit's declared margin of the matched random-network mean in \nm{ben13.for.random-level} of the \nm{ben13.cells} cells. The random-network comparison was unavailable in the other cell.

\paragraph{Recorded runs.} The instrumented Causal-Copilot series records code identity, the model snapshot,
model requests and replies, and stage graphs and refinement steps as captured by the run.
Causal-Copilot was run unmodified (upstream commit \texttt{a31aa7c}) with \nm{bnd08.model}: \nm{bnd08.runs} runs on the
semantic suite in the name-arm study, and a main-control series of \nm{cc.arms} runs on the same tables (API cost
\nm{cc.usd}). The other systems' name-control results use a different semantic cohort and their own declared output/scoring
contracts. Their results are not paired scenario-for-scenario with the Causal-Copilot collection.

\paragraph{Limits of the data.} The semantic scenarios are simulated, so their exactness comes with simulation's
artefacts; each table and arm was run once; one model family was used; and the named datasets' references are expert
graphs. The protocol treats each of these as a declared regime or tier rather than hiding it\ours.

\section*{Code}
The reusable framework is published at \url{https://github.com/crucis-lab/crucis}, with documentation at
\url{https://crucis-lab.github.io/crucis/}.
The companion repository (\url{https://github.com/uzairgheewala/causal-copilot-q1}) uses retained graphs to
recompute its Causal-Copilot run comparison with one command and no API key. The stage-specific attribution uses
separate evidence in the research repository (\url{https://github.com/uzairgheewala/cc-research}).
Appendix~\ref{app:numbers} maps reported values to result files, retained diagnostic snapshots and separately
preserved CausalBench metadata. The numerical evidence is pinned to \texttt{\nm{cc.head}}.

%%%%%%%%%%%%%%%%%%%%%%%%%%%%%%%%%%%%%%%%%%%%%%%%%%%%%%%%
%%%% Reference / Bibliography
%%%%%%%%%%%%%%%%%%%%%%%%%%%%%%%%%%%%%%%%%%%%%%%%%%%%%%%%

\makereference

\begingroup
\checkpointbibstyle
\bibliography{reference}
\bibliographystyle{style/dsc180bibstyle}
\endgroup

%%%%%%%%%%%%%%%%%%%%%%%%%%%%%%%%%%%%%%%%%%%%%%%%%%%%%%%%
%%%% Appendix
%%%%%%%%%%%%%%%%%%%%%%%%%%%%%%%%%%%%%%%%%%%%%%%%%%%%%%%%

\clearpage
\makeappendix

Use this appendix to follow an audit from its question to the evidence behind the
answer. Start with scoring (Appendix~\ref{app:scoring}), names
(Appendix~\ref{app:names}), or pipeline stages (Appendix~\ref{app:stages}).
Appendix~\ref{app:numbers} is the lookup for exact values and their sources.

\subsection{Scoring the same graphs}
\label{app:scoring}
CausalArena's released files list the variables in causal order. The scoring rule
uses column position to settle directions the algorithm left uncertain. On the evaluated semantic subset, FCI gains
\nm{ab01.fci.gain} directed F1 under that rule relative to expectation scoring on the same graphs.
The published PC row is \nm{ab01.pc.published}; a distinct PC maxP variant scores \nm{ab01.pcmaxp.controlled}
on shuffled columns under expectation scoring. The FCI and maxP rows cover subsets
of the public inventory described in Section~\ref{sec:data}.

The separate CausalBench comparison in Section~\ref{sec:data} asks what the metrics
reward when a reference graph is unavailable. Its smaller edge budget, test-informed
reference and unavailable random-network comparison limit what can be inferred.

\subsection{Changing names while holding the data fixed}
\label{app:names}
The name control pairs graphs produced from the same numerical table under real
and swapped variable names. The AB-06 tally applies a \nm{ab06.p9.min-scen}-pair
entry floor to each system's declared output stage. PyWhy-LLM falls below that floor.

Swapping names lowers directed F1 by at least \nm{ab06.p9.threshold} in
\nm{ab06.p9.losing} of \nm{ab06.p9.entered} entries. The drops are
\nm{ab06.causaliq.drop} for causaliq-knowledge, \nm{ab06.cc.drop} for Causal-Copilot,
and \nm{ab06.matmcd.drop} for MATMCD. The pipeline of \citet{vashishtha2025causal}
shows a drop of \nm{ab06.vashishtha.drop} under its tested conditions.

The tally describes sensitivity under each entry's measurement. The earlier
exploratory Causal-Copilot comparison scores committed marks in the graph handed
to cycle removal. The other agents use final outputs and expectation scoring on
a semantic cohort with no shared scenarios. The tally therefore cannot rank
these systems under common conditions.

\paragraph{Which reading do the retained changes follow?}
Inside Causal-Copilot, BND-08's recorded label-following calculation counts distinct
(scenario, directed-edge) pairs retained in completed swapped-name outputs, changed from bootstrap and matching
their post-review pair state at final output, and adjacent in exactly one reference reading.
Among these \nm{bnd08.m3.n} informative pairs, \nm{bnd08.m3.labels} are adjacent under the label reading and
\nm{bnd08.m3.data} under the data reading. This classification tests adjacency, not agreement of edge direction;
edges adjacent in both references or neither are outside this denominator. The recorded decision rule comes from BND-08's
pre-run specification, not a later finding threshold. However, that specification describes a pre-cycle endpoint
while the recorded calculation consumes completed final outputs. We retain the counts as a descriptive diagnostic,
without claiming confirmation of the specified endpoint.
Failed swapped runs cannot enter this completed-output calculation. Edges from
the same scenario are dependent. The preserved manifest binds the rule's bytes.
It does not establish when the rule became public or reconcile the different
endpoints.

\subsection{Locating a change inside a pipeline}
\label{app:stages}
Causal-Copilot's main-control initial mean adjacency F1 is \nm{cc.s0-adj} across all \nm{cc.arms} tables.
Its final mean is \nm{cc.final-adj} among the \nm{cc.completed} completions. In paired diagnostics on those
completions, the largest observed stage loss occurs at bootstrap within refinement, before model review.
The initial and final means above describe different populations. The stage
comparison instead follows the same completed runs through successive graphs.
A real-arm-only effect answers a different question and is not substituted here.

The completed-run comparison leaves out a separate failure mode. All
\nm{cc.cyc.entered} runs that reached cycle removal crashed there. The logs
record a key error in \nm{cc.cyc.keyerr} of those entrants.

\begin{figure}[htbp]
\centering
\small
\fbox{Initial graph}
\quad\textgreater\quad
\colorbox{StageHighlight}{\textbf{Bootstrap refinement}}
\quad\textgreater\quad
\fbox{Language-model review}
\caption{The bootstrap transition had the largest paired drop in adjacency F1
among the completed Causal-Copilot main-control runs. The shaded stage
locates that transition before language-model review.}
\label{fig:stage-locator}
\end{figure}

\paragraph{A separate comparison against a data-only graph.}
In the real-name MATMCD runs,
pooled precision among committed claims on the released intervention targets is lower for the final output than
for the data-only variant on the point estimate (\nm{ab06.matmcd.dv1} against \nm{ab06.matmcd.dv1-data}); the
interval for the difference includes zero. The pre-registered point rule
applies to the estimate.

\subsection{Where every number comes from}
\label{app:numbers}
Generated by \texttt{docs/paper/q1/build\_numbers.py} from \texttt{numbers.json}. Paths are relative to the root of
cc-research. Selectors extend JSON pointers: \texttt{*} takes every child and \texttt{[k=v]} selects the first
matching list item. Source flags identify inputs outside the tracked build base or with differences from it.
The CausalBench filtering metadata are separately preserved and hashed. The table locates evidence for each value;
the companion replay covers only its stated comparison.

\begingroup
\let\small\footnotesize
\renewcommand{\_}{\textunderscore\allowbreak}
% Generated by build_numbers.py. Do not edit.
{\small\begin{longtable}{p{0.27\linewidth} p{0.11\linewidth} p{0.54\linewidth}}
\toprule Key & Value & Source file and JSON pointer (cc-research) \\ \midrule \endhead
\texttt{cc.arms} & 30 & {\footnotesize \texttt{analysis/\allowbreak{}results/\allowbreak{}fork/\allowbreak{}bld06\_live\_compare.json} \texttt{/\allowbreak{}arms}} \\
\texttt{cc.completed} & 22 & {\footnotesize \texttt{analysis/\allowbreak{}results/\allowbreak{}fork/\allowbreak{}bld06\_live\_compare.json} \texttt{/\allowbreak{}completed/\allowbreak{}main}} \\
\texttt{cc.s0-adj} & 0.584 & {\footnotesize \texttt{analysis/\allowbreak{}results/\allowbreak{}fork/\allowbreak{}bld06\_live\_compare.json} \texttt{/\allowbreak{}means/\allowbreak{}main/\allowbreak{}S0\_adj\_f1}} \\
\texttt{cc.final-adj} & 0.321 & {\footnotesize \texttt{analysis/\allowbreak{}results/\allowbreak{}fork/\allowbreak{}bld06\_live\_compare.json} \texttt{/\allowbreak{}means/\allowbreak{}main/\allowbreak{}final\_adj\_f1}} \\
\texttt{cc.usd} & \$12.57 & {\footnotesize \texttt{analysis/\allowbreak{}results/\allowbreak{}fork/\allowbreak{}bld06\_live\_compare.json} \texttt{/\allowbreak{}usd/\allowbreak{}main}} \\
\texttt{cc.cyc.entered} & 8 & {\footnotesize \texttt{analysis/\allowbreak{}results/\allowbreak{}fork/\allowbreak{}bld06\_cycles.json} \texttt{/\allowbreak{}arms/\allowbreak{}main/\allowbreak{}runs\_entering\_cycle\_step}} \\
\texttt{cc.cyc.keyerr} & 8 & {\footnotesize \texttt{analysis/\allowbreak{}results/\allowbreak{}fork/\allowbreak{}bld06\_cycles.json} \texttt{/\allowbreak{}arms/\allowbreak{}main/\allowbreak{}cycle\_step\_keyerrors}} \\
\texttt{bnd08.runs} & 30 & {\footnotesize len of: \texttt{analysis/\allowbreak{}results/\allowbreak{}stability\_s2/\allowbreak{}bnd08\_summary.json} \texttt{/\allowbreak{}runs}} \\
\texttt{bnd08.m3.labels} & 81 & {\footnotesize \texttt{analysis/\allowbreak{}results/\allowbreak{}stability\_s2/\allowbreak{}bnd08\_summary.json} \texttt{/\allowbreak{}M3/\allowbreak{}labels\_only}} \\
\texttt{bnd08.m3.data} & 20 & {\footnotesize \texttt{analysis/\allowbreak{}results/\allowbreak{}stability\_s2/\allowbreak{}bnd08\_summary.json} \texttt{/\allowbreak{}M3/\allowbreak{}data\_only}} \\
\texttt{bnd08.m3.n} & 101 & {\footnotesize \texttt{analysis/\allowbreak{}results/\allowbreak{}stability\_s2/\allowbreak{}bnd08\_summary.json} \texttt{/\allowbreak{}M3/\allowbreak{}n}} \\
\texttt{bnd08.model} & gpt-4o-2024-08-06 & {\footnotesize \texttt{analysis/\allowbreak{}results/\allowbreak{}stability\_s2/\allowbreak{}bnd08\_summary.json} \texttt{/\allowbreak{}snapshots/\allowbreak{}0}} \\
\texttt{ab06.agents} & 4 & {\footnotesize len of: \texttt{analysis/\allowbreak{}results/\allowbreak{}ab06/\allowbreak{}ab06.json} \texttt{/\allowbreak{}agents}} \\
\texttt{ab06.cc-entered} & 1 & {\footnotesize count\_true of: \texttt{analysis/\allowbreak{}results/\allowbreak{}ab06/\allowbreak{}ab06.json} \texttt{/\allowbreak{}causal\_copilot/\allowbreak{}*/\allowbreak{}counted\_in\_P9}} \\
\texttt{ab06.systems} & 5 & {\footnotesize add of: key \texttt{ab06.agents}; key \texttt{ab06.cc-entered}} \\
\texttt{ab06.p9.entered} & 4 & {\footnotesize len of: \texttt{analysis/\allowbreak{}results/\allowbreak{}ab06/\allowbreak{}ab06.json} \texttt{/\allowbreak{}P9/\allowbreak{}systems\_entered}} \\
\texttt{ab06.p9.losing} & 3 & {\footnotesize len of: \texttt{analysis/\allowbreak{}results/\allowbreak{}ab06/\allowbreak{}ab06.json} \texttt{/\allowbreak{}P9/\allowbreak{}losing\_ge\_0.10}} \\
\texttt{ab06.p9.min-scen} & 5 & {\footnotesize \texttt{analysis/\allowbreak{}results/\allowbreak{}ab06/\allowbreak{}ab06.json} \texttt{/\allowbreak{}P9/\allowbreak{}min\_paired\_scenarios}} \\
\texttt{ab06.p9.threshold} & 0.10 & {\footnotesize \texttt{plans/\allowbreak{}ab06\_prereg.json} \texttt{/\allowbreak{}predictions/\allowbreak{}P9/\allowbreak{}supported} (pattern \texttt{lose >= ([0-9.]+)})} \\
\texttt{ab06.causaliq.drop} & 0.251 & {\footnotesize \texttt{analysis/\allowbreak{}results/\allowbreak{}ab06/\allowbreak{}ab06.json} \texttt{/\allowbreak{}agents/\allowbreak{}causaliq/\allowbreak{}name\_reliance\_real\_minus\_swapped/\allowbreak{}mean}} \\
\texttt{ab06.cc.drop} & 0.167 & {\footnotesize \texttt{analysis/\allowbreak{}results/\allowbreak{}ab06/\allowbreak{}ab06.json} \texttt{/\allowbreak{}causal\_copilot/\allowbreak{}causal\_copilot\_gpt4o/\allowbreak{}name\_reliance\_real\_minus\_swapped/\allowbreak{}mean}} \\
\texttt{ab06.matmcd.drop} & 0.122 & {\footnotesize \texttt{analysis/\allowbreak{}results/\allowbreak{}ab06/\allowbreak{}ab06.json} \texttt{/\allowbreak{}agents/\allowbreak{}matmcd/\allowbreak{}name\_reliance\_real\_minus\_swapped/\allowbreak{}mean}} \\
\texttt{ab06.vashishtha.drop} & 0.020 & {\footnotesize \texttt{analysis/\allowbreak{}results/\allowbreak{}ab06/\allowbreak{}ab06.json} \texttt{/\allowbreak{}agents/\allowbreak{}vashishtha/\allowbreak{}name\_reliance\_real\_minus\_swapped/\allowbreak{}mean}} \\
\texttt{ab06.matmcd.dv1} & 0.6855 & {\footnotesize \texttt{analysis/\allowbreak{}results/\allowbreak{}ab06/\allowbreak{}ab06.json} \texttt{/\allowbreak{}agents/\allowbreak{}matmcd/\allowbreak{}data\_only\_floor\_real\_arm/\allowbreak{}dv1\_precision\_final}} \\
\texttt{ab06.matmcd.dv1-data} & 0.7199 & {\footnotesize \texttt{analysis/\allowbreak{}results/\allowbreak{}ab06/\allowbreak{}ab06.json} \texttt{/\allowbreak{}agents/\allowbreak{}matmcd/\allowbreak{}data\_only\_floor\_real\_arm/\allowbreak{}dv1\_precision\_data\_only}} \\
\texttt{ab01.pc.published} & 0.333 & {\footnotesize \texttt{analysis/\allowbreak{}results/\allowbreak{}rescore\_v2/\allowbreak{}leaderboard\_controlled.json} \texttt{/\allowbreak{}table/\allowbreak{}[method=PC (upstream, default rule);family=semantic]/\allowbreak{}published}} \\
\texttt{ab01.pcmaxp.controlled} & 0.463 & {\footnotesize \texttt{analysis/\allowbreak{}results/\allowbreak{}rescore\_v2/\allowbreak{}leaderboard\_controlled.json} \texttt{/\allowbreak{}table/\allowbreak{}[method=PC (upstream, maxP);family=semantic]/\allowbreak{}ours\_shuffled\_expectation}} \\
\texttt{ab01.fci.authors} & 0.578 & {\footnotesize \texttt{analysis/\allowbreak{}results/\allowbreak{}rescore\_v2/\allowbreak{}leaderboard\_controlled.json} \texttt{/\allowbreak{}table/\allowbreak{}[method=FCI;family=semantic]/\allowbreak{}ours\_released\_authors\_rule}} \\
\texttt{ab01.fci.expect} & 0.357 & {\footnotesize \texttt{analysis/\allowbreak{}results/\allowbreak{}rescore\_v2/\allowbreak{}leaderboard\_controlled.json} \texttt{/\allowbreak{}table/\allowbreak{}[method=FCI;family=semantic]/\allowbreak{}ours\_released\_expectation}} \\
\texttt{ab01.fci.gain} & 0.221 & {\footnotesize sub of: key \texttt{ab01.fci.authors}; key \texttt{ab01.fci.expect}} \\
\texttt{ab01.files.semantic} & 250 & {\footnotesize \texttt{analysis/\allowbreak{}results/\allowbreak{}rescore\_v2/\allowbreak{}leaderboard\_controlled.json} \texttt{/\allowbreak{}table/\allowbreak{}[method=DirectLiNGAM;family=semantic]/\allowbreak{}files}} \\
\texttt{ab01.files.formula} & 250 & {\footnotesize \texttt{analysis/\allowbreak{}results/\allowbreak{}rescore\_v2/\allowbreak{}leaderboard\_controlled.json} \texttt{/\allowbreak{}table/\allowbreak{}[method=DirectLiNGAM;family=formula]/\allowbreak{}files}} \\
\texttt{ab01.files.synthetic} & 100 & {\footnotesize \texttt{analysis/\allowbreak{}results/\allowbreak{}rescore\_v2/\allowbreak{}leaderboard\_controlled.json} \texttt{/\allowbreak{}table/\allowbreak{}[method=DirectLiNGAM;family=synthetic]/\allowbreak{}files}} \\
\texttt{ben13.reproduced} & 4 & {\footnotesize \texttt{analysis/\allowbreak{}results/\allowbreak{}ben13/\allowbreak{}summary.json} \texttt{/\allowbreak{}reproduced\_rows}} \\
\texttt{ben13.cells} & 4 & {\footnotesize count\_notnull of: \texttt{analysis/\allowbreak{}results/\allowbreak{}ben13/\allowbreak{}summary.json} \texttt{/\allowbreak{}lines/\allowbreak{}*/\allowbreak{}controls/\allowbreak{}*/\allowbreak{}placement\_share\_wasserstein}} \\
\texttt{ben13.for.random-level} & 3 & {\footnotesize count\_notnull of: \texttt{analysis/\allowbreak{}results/\allowbreak{}ben13/\allowbreak{}summary.json} \texttt{/\allowbreak{}lines/\allowbreak{}*/\allowbreak{}controls/\allowbreak{}*/\allowbreak{}?placement\_share\_for\_note}} \\
\texttt{ben13.echo.min} & 0.90 & {\footnotesize min of: \texttt{analysis/\allowbreak{}results/\allowbreak{}ben13/\allowbreak{}summary.json} \texttt{/\allowbreak{}predictions/\allowbreak{}P3/\allowbreak{}values/\allowbreak{}*}} \\
\texttt{ben13.echo.max} & 0.95 & {\footnotesize max of: \texttt{analysis/\allowbreak{}results/\allowbreak{}ben13/\allowbreak{}summary.json} \texttt{/\allowbreak{}predictions/\allowbreak{}P3/\allowbreak{}values/\allowbreak{}*}} \\
\texttt{s2.families} & 10 & {\footnotesize len of: \texttt{datasets/\allowbreak{}suite\_q1s2/\allowbreak{}MAPPING.json} \texttt{/\allowbreak{}families}} \\
\texttt{s2.tables} & 30 & {\footnotesize len of: \texttt{datasets/\allowbreak{}suite\_q1s2/\allowbreak{}MAPPING.json} \texttt{/\allowbreak{}arms}} \\
\texttt{s2.arms-per} & 3 & {\footnotesize div of: key \texttt{s2.tables}; key \texttt{s2.families}} \\
\texttt{s2.rows} & 1,000 & {\footnotesize max of: \texttt{datasets/\allowbreak{}suite\_q1s2/\allowbreak{}MAPPING.json} \texttt{/\allowbreak{}families/\allowbreak{}*/\allowbreak{}rows}} \\
\texttt{s2.vars.min} & 20 & {\footnotesize min of: \texttt{datasets/\allowbreak{}suite\_q1s2/\allowbreak{}MAPPING.json} \texttt{/\allowbreak{}families/\allowbreak{}*/\allowbreak{}variables}} \\
\texttt{s2.vars.max} & 25 & {\footnotesize max of: \texttt{datasets/\allowbreak{}suite\_q1s2/\allowbreak{}MAPPING.json} \texttt{/\allowbreak{}families/\allowbreak{}*/\allowbreak{}variables}} \\
\texttt{s2.edges.min} & 35 & {\footnotesize min of: \texttt{datasets/\allowbreak{}suite\_q1s2/\allowbreak{}MAPPING.json} \texttt{/\allowbreak{}families/\allowbreak{}*/\allowbreak{}edges}} \\
\texttt{s2.edges.max} & 58 & {\footnotesize max of: \texttt{datasets/\allowbreak{}suite\_q1s2/\allowbreak{}MAPPING.json} \texttt{/\allowbreak{}families/\allowbreak{}*/\allowbreak{}edges}} \\
\texttt{s1.sachs.rows} & 853 & {\footnotesize \texttt{datasets/\allowbreak{}MAPPING.json} \texttt{/\allowbreak{}source/\allowbreak{}rows}} \\
\texttt{s1.sachs.vars} & 11 & {\footnotesize len of: \texttt{datasets/\allowbreak{}MAPPING.json} \texttt{/\allowbreak{}source/\allowbreak{}columns}} \\
\texttt{s1.sachs.edges} & 19 & {\footnotesize len of: \texttt{datasets/\allowbreak{}MAPPING.json} \texttt{/\allowbreak{}source/\allowbreak{}reference\_edges}} \\
\texttt{s1.sachs.arms} & 4 & {\footnotesize len of: \texttt{datasets/\allowbreak{}MAPPING.json} \texttt{/\allowbreak{}arms}} \\
\texttt{s1.sachs.bnlearn} & 17 & {\footnotesize \texttt{analysis/\allowbreak{}results/\allowbreak{}reference\_reliability/\allowbreak{}reliability.json} \texttt{/\allowbreak{}families/\allowbreak{}sachs/\allowbreak{}repo\_vs\_bnlearn/\allowbreak{}edges\_b}} \\
\texttt{s1.sachs.shared} & 16 & {\footnotesize \texttt{analysis/\allowbreak{}results/\allowbreak{}reference\_reliability/\allowbreak{}reliability.json} \texttt{/\allowbreak{}families/\allowbreak{}sachs/\allowbreak{}repo\_vs\_bnlearn/\allowbreak{}shared\_adjacencies}} \\
\texttt{s1.sachs.reversed} & 1 & {\footnotesize \texttt{analysis/\allowbreak{}results/\allowbreak{}reference\_reliability/\allowbreak{}reliability.json} \texttt{/\allowbreak{}families/\allowbreak{}sachs/\allowbreak{}repo\_vs\_bnlearn/\allowbreak{}reversed}} \\
\texttt{s1.sachs.arena} & 20 & {\footnotesize \texttt{analysis/\allowbreak{}results/\allowbreak{}reference\_reliability/\allowbreak{}reliability.json} \texttt{/\allowbreak{}families/\allowbreak{}sachs/\allowbreak{}causalarena\_version/\allowbreak{}edges}} \\
\texttt{s1.abalone.rows} & 4,177 & {\footnotesize \texttt{datasets/\allowbreak{}suite\_b1/\allowbreak{}MAPPING.json} \texttt{/\allowbreak{}families/\allowbreak{}abalone/\allowbreak{}rows}} \\
\texttt{s1.abalone.vars} & 8 & {\footnotesize len of: \texttt{datasets/\allowbreak{}suite\_b1/\allowbreak{}MAPPING.json} \texttt{/\allowbreak{}families/\allowbreak{}abalone/\allowbreak{}columns}} \\
\texttt{s1.abalone.edges} & 7 & {\footnotesize len of: \texttt{datasets/\allowbreak{}suite\_b1/\allowbreak{}MAPPING.json} \texttt{/\allowbreak{}families/\allowbreak{}abalone/\allowbreak{}reference\_edges}} \\
\texttt{s1.autompg.rows} & 392 & {\footnotesize \texttt{datasets/\allowbreak{}suite\_b1/\allowbreak{}MAPPING.json} \texttt{/\allowbreak{}families/\allowbreak{}autompg/\allowbreak{}rows}} \\
\texttt{s1.autompg.vars} & 5 & {\footnotesize len of: \texttt{datasets/\allowbreak{}suite\_b1/\allowbreak{}MAPPING.json} \texttt{/\allowbreak{}families/\allowbreak{}autompg/\allowbreak{}columns}} \\
\texttt{s1.autompg.edges} & 5 & {\footnotesize len of: \texttt{datasets/\allowbreak{}suite\_b1/\allowbreak{}MAPPING.json} \texttt{/\allowbreak{}families/\allowbreak{}autompg/\allowbreak{}reference\_edges}} \\
\texttt{s1.concrete.rows} & 1,030 & {\footnotesize \texttt{datasets/\allowbreak{}suite\_b1/\allowbreak{}MAPPING.json} \texttt{/\allowbreak{}families/\allowbreak{}concrete/\allowbreak{}rows}} \\
\texttt{s1.concrete.vars} & 9 & {\footnotesize len of: \texttt{datasets/\allowbreak{}suite\_b1/\allowbreak{}MAPPING.json} \texttt{/\allowbreak{}families/\allowbreak{}concrete/\allowbreak{}columns}} \\
\texttt{s1.concrete.edges} & 8 & {\footnotesize len of: \texttt{datasets/\allowbreak{}suite\_b1/\allowbreak{}MAPPING.json} \texttt{/\allowbreak{}families/\allowbreak{}concrete/\allowbreak{}reference\_edges}} \\
\texttt{s1.dwd.rows} & 349 & {\footnotesize \texttt{datasets/\allowbreak{}suite\_b1/\allowbreak{}MAPPING.json} \texttt{/\allowbreak{}families/\allowbreak{}dwd/\allowbreak{}rows}} \\
\texttt{s1.dwd.vars} & 5 & {\footnotesize len of: \texttt{datasets/\allowbreak{}suite\_b1/\allowbreak{}MAPPING.json} \texttt{/\allowbreak{}families/\allowbreak{}dwd/\allowbreak{}columns}} \\
\texttt{s1.dwd.edges} & 4 & {\footnotesize len of: \texttt{datasets/\allowbreak{}suite\_b1/\allowbreak{}MAPPING.json} \texttt{/\allowbreak{}families/\allowbreak{}dwd/\allowbreak{}reference\_edges}} \\
\texttt{s1.autompg.arena.vars} & 4 & {\footnotesize \texttt{analysis/\allowbreak{}results/\allowbreak{}reference\_reliability/\allowbreak{}reliability.json} \texttt{/\allowbreak{}families/\allowbreak{}autompg/\allowbreak{}causalarena\_vars\_edges\_rows/\allowbreak{}0}} \\
\texttt{s1.autompg.arena.edges} & 3 & {\footnotesize \texttt{analysis/\allowbreak{}results/\allowbreak{}reference\_reliability/\allowbreak{}reliability.json} \texttt{/\allowbreak{}families/\allowbreak{}autompg/\allowbreak{}causalarena\_vars\_edges\_rows/\allowbreak{}1}} \\
\texttt{s1.autompg.arena.rows} & 398 & {\footnotesize \texttt{analysis/\allowbreak{}results/\allowbreak{}reference\_reliability/\allowbreak{}reliability.json} \texttt{/\allowbreak{}families/\allowbreak{}autompg/\allowbreak{}causalarena\_vars\_edges\_rows/\allowbreak{}2}} \\
\texttt{s1.named.tables} & 12 & {\footnotesize len of: \texttt{datasets/\allowbreak{}suite\_b1/\allowbreak{}MAPPING.json} \texttt{/\allowbreak{}arms}} \\
\texttt{s1.named} & 4 & {\footnotesize len of: \texttt{datasets/\allowbreak{}suite\_b1/\allowbreak{}MAPPING.json} \texttt{/\allowbreak{}families}} \\
\texttt{s1.named.arms-per} & 3 & {\footnotesize div of: key \texttt{s1.named.tables}; key \texttt{s1.named}} \\
\texttt{cb.k562.cells} & 162,751 & {\footnotesize \texttt{external/\allowbreak{}causalbench/\allowbreak{}dataset\_k562\_filtered.npz.json} \texttt{/\allowbreak{}cells\_out} \emph{(excluded)}} \\
\texttt{cb.k562.genes} & 622 & {\footnotesize \texttt{external/\allowbreak{}causalbench/\allowbreak{}dataset\_k562\_filtered.npz.json} \texttt{/\allowbreak{}genes\_out} \emph{(excluded)}} \\
\texttt{cb.k562.controls} & 10,691 & {\footnotesize \texttt{external/\allowbreak{}causalbench/\allowbreak{}dataset\_k562\_filtered.npz.json} \texttt{/\allowbreak{}controls} \emph{(excluded)}} \\
\texttt{cb.rpe1.cells} & 162,733 & {\footnotesize \texttt{external/\allowbreak{}causalbench/\allowbreak{}dataset\_rpe1\_filtered.npz.json} \texttt{/\allowbreak{}cells\_out} \emph{(excluded)}} \\
\texttt{cb.rpe1.genes} & 383 & {\footnotesize \texttt{external/\allowbreak{}causalbench/\allowbreak{}dataset\_rpe1\_filtered.npz.json} \texttt{/\allowbreak{}genes\_out} \emph{(excluded)}} \\
\texttt{cb.rpe1.controls} & 11,485 & {\footnotesize \texttt{external/\allowbreak{}causalbench/\allowbreak{}dataset\_rpe1\_filtered.npz.json} \texttt{/\allowbreak{}controls} \emph{(excluded)}} \\
\bottomrule\end{longtable}}

\endgroup

\end{document}
